\chapter{Sprint 3 - Visual AI Module \& Voice Agents (LiveKit)}
\label{chap:sprint3}

\section*{Introduction}

\noindent This chapter presents the third development sprint, which adds multimodal support capabilities to the platform. Sprint 3 focuses on two complementary extensions: visual understanding of user screens and real-time voice assistance through LiveKit AI agents. The implemented source code turns screenshots, screen recordings, and support calls into structured support context that can be reused by the ticketing, RAG, and escalation workflows.

\medskip

\textbf{Sprint Goal:} Implement a Visual AI module for screenshot and screenshare analysis, and deploy LiveKit voice agents capable of answering customer questions, reading live shared-screen context, querying the RAG knowledge base, recording calls, and escalating unresolved cases to human support.

%---------------------------------------------------------------------
\section{Sprint Backlog}

This backlog reflects the implemented Sprint 3 scope in the current source code. It includes the Visual AI persistence model, Gemini-based visual analysis, screenshare assistance, LiveKit agent orchestration, internal service-to-service APIs, and voice call logging.

\begin{table}[H]
\smaller
\centering
\caption{Sprint 3 Backlog}
\label{tab:sprint3-backlog}
\begin{tabular}{|l|p{4.8cm}|p{4.8cm}|c|}
\hline
\textbf{US \#} & \textbf{Feature} & \textbf{Task} & \textbf{Est.\,(h)} \\
\hline
\multirow{4}{*}{US4} & \multirow{4}{*}{\parbox{4.8cm}{VISUAL AI PIPELINE}} & Screenshot upload, consent validation, and filesystem storage & 15 \\
\cline{3-4}
 & & Gemini Vision OCR, UI captioning, element detection, and embeddings with CLIP fallback & 35 \\
\cline{3-4}
 & & Reference screen management and UI gap detection & 30 \\
\cline{3-4}
 & & UI state timeline for conversation-linked screenshots & 20 \\
\hline
\multirow{3}{*}{US9} & \multirow{3}{*}{\parbox{4.8cm}{SCREENSHARE ASSISTANCE}} & Low-FPS frame sampling and transition scoring & 20 \\
\cline{3-4}
 & & Uploaded video and realtime chunk processing with FFmpeg & 25 \\
\cline{3-4}
 & & Support-call screen context bridge for voice agents & 25 \\
\hline
\multirow{3}{*}{US13} & \multirow{3}{*}{\parbox{4.8cm}{LIVEKIT VOICE AGENTS}} & LiveKit server and voice-agent worker container & 30 \\
\cline{3-4}
 & & StarterAgent ``Tom'' and specialist routing tools & 25 \\
\cline{3-4}
 & & Support, Booking, and FAQ agents with multilingual prompts & 35 \\
\hline
\multirow{3}{*}{US14} & \multirow{3}{*}{\parbox{4.8cm}{VOICE OPERATIONS}} & Internal RAG bridge and service-key authentication & 20 \\
\cline{3-4}
 & & STT/TTS/LLM factory: Gemini, OpenAI, local Whisper, and Edge TTS & 25 \\
\cline{3-4}
 & & Call transcript capture, audio recording, post-call summary, and ticket linking & 30 \\
\hline
\multicolumn{3}{|r|}{\textbf{Total}} & \textbf{260} \\
\hline
\end{tabular}
\end{table}

%---------------------------------------------------------------------
\section{Design}

\subsection{Database Schema}

Sprint 3 does not use generic tables such as \texttt{visual\_contexts} or \texttt{agent\_sessions}. The implemented schema stores visual artifacts, analyses, reference screens, UI timeline states, and LiveKit call logs in dedicated tables.

\begin{table}[H]
\centering
\small
\caption{Sprint 3 - Multimodal Database Tables}
\label{tab:sprint3-db-tables}
\begin{tabular}{|l|p{9.5cm}|}
\hline
\textbf{Table} & \textbf{Key Columns} \\
\hline
\texttt{screenshots} & \texttt{id}, \texttt{conversation\_id}, \texttt{user\_id}, \texttt{filename}, \texttt{file\_path}, \texttt{file\_size}, \texttt{mime\_type}, \texttt{consent}, \texttt{metadata}, timestamps, soft-delete fields \\
\hline
\texttt{visual\_analyses} & \texttt{id}, \texttt{screenshot\_id}, \texttt{provider}, \texttt{ocr\_text}, \texttt{caption}, \texttt{elements}, \texttt{labels}, \texttt{regions}, \texttt{embedding}, \texttt{raw\_result}, \texttt{confidence}, \texttt{processing\_ms} \\
\hline
\texttt{ui\_states} & \texttt{id}, \texttt{conversation\_id}, \texttt{analysis\_id}, \texttt{screenshot\_id}, \texttt{state\_label}, \texttt{state\_data}, \texttt{embedding}, \texttt{sequence\_num}, \texttt{gap\_detected}, \texttt{gap\_details}, \texttt{gap\_severity} \\
\hline
\texttt{reference\_screens} & \texttt{id}, \texttt{name}, \texttt{description}, \texttt{screen\_key}, \texttt{file\_path}, \texttt{embedding}, \texttt{expected\_elements}, \texttt{expected\_ocr\_text} \\
\hline
\texttt{voice\_call\_logs} & \texttt{id}, \texttt{room\_name}, \texttt{room\_sid}, \texttt{transcript}, \texttt{audio\_file\_path}, \texttt{duration\_seconds}, \texttt{started\_at}, \texttt{ended\_at}, timestamps \\
\hline
\texttt{tickets} & Includes \texttt{source\_voice\_call\_id} so a call log can be attached to a new or existing support ticket \\
\hline
\end{tabular}
\end{table}

\subsection{Sequence Diagrams}

The Visual AI workflow analyzes either a stored screenshot or a live screenshare packet. In the support-call path, the frontend captures frames or short video chunks, sends them to the Visual AI API, and publishes the resulting text/context to a short-lived support-call context store. Voice agents then read that context through the internal API.

\begin{figure}[H]
\centering
\includegraphics[width=0.93\textwidth]{sequence-diagram-screensharing.png}
\caption{Sequence Diagram - Visual AI Screen Sharing and Voice-Agent Context Bridge}
\label{fig:sprint3-screensharing}
\end{figure}

%---------------------------------------------------------------------
\section{Implementation}

\subsection{Visual AI Module}

The Visual AI module is implemented under \texttt{app/visual\_ai}. Its public API is mounted at \texttt{/api/v1/visual-ai}. The current provider registry forces the Gemini provider, while retaining compatibility with older provider parameters. The Gemini provider performs OCR extraction, UI captioning, element detection, and visual embedding generation. If Gemini embeddings are unavailable, the implementation falls back to CLIP image embeddings.

\begin{itemize}
    \item \textbf{Screenshot pipeline:} \texttt{POST /visual-ai/upload} stores consented images on disk and creates a \texttt{Screenshot} row. \texttt{POST /visual-ai/screenshots/\{id\}/analyze} creates a \texttt{VisualAnalysis} row containing OCR text, caption, UI elements, labels, regions, embedding, confidence, and processing duration.
    \item \textbf{Raw and full analysis:} \texttt{POST /visual-ai/analyze-raw} analyzes an image without storing it. \texttt{POST /visual-ai/process} combines upload, analysis, optional reference gap detection, timeline insertion, and guidance generation.
    \item \textbf{Reference screens:} Admins create reference screens through \texttt{POST /visual-ai/references}. Each reference has a unique \texttt{screen\_key}, an image path, optional expected OCR text, optional expected UI elements, and an embedding.
    \item \textbf{Gap detection:} The gap detector combines visual similarity, OCR text difference, UI element difference, and visible error/loading penalties. Scores map to \texttt{NO\_GAP}, \texttt{MINOR}, \texttt{SIGNIFICANT}, or \texttt{CRITICAL}.
    \item \textbf{Timeline tracking:} When a screenshot is linked to a conversation, the module creates ordered \texttt{ui\_states} entries. This gives agents a chronological view of what the user saw during troubleshooting.
    \item \textbf{Screenshare assistance:} \texttt{/visual-ai/screenshare/assist} accepts ordered frame images. \texttt{/visual-ai/screenshare/assist-video} extracts frames from an uploaded video. \texttt{/visual-ai/screenshare/assist-realtime-chunk} processes short realtime recording chunks. All three return transition scores, final-frame OCR/caption data, embedding backend details, optional reference similarity, and assistance hints.
    \item \textbf{Troubleshooting wizard:} \texttt{POST /visual-ai/troubleshooting/wizard} generates deterministic step-by-step troubleshooting guidance from the user goal, observed caption, visible text, attempted actions, and context hints.
\end{itemize}

\subsection{Real-time Voice Agents (LiveKit)}

The LiveKit voice subsystem lives under \texttt{voice\_agents}. It runs as a separate Docker service and exposes a small control API on port \texttt{8001}. The backend proxies this control API through \texttt{/api/v1/voice-agents}, allowing admins to start, stop, inspect logs, update runtime configuration, and generate LiveKit tokens.

\begin{itemize}
    \item \textbf{Agent roster:} \texttt{StarterAgent} (Tom) greets users and routes them to \texttt{SupportAgent} (Mike), \texttt{BookingAgent} (Jessica), or \texttt{FAQAgent} (Ameni). Agents can return to Tom when needed.
    \item \textbf{Multilingual behavior:} Agent prompts support English, French, Modern Standard Arabic, and Tunisian Derja. They answer in the user's language and include empathy and de-escalation instructions.
    \item \textbf{RAG grounding:} All agents inherit \texttt{search\_knowledge\_base} and \texttt{generate\_answer} tools. These call \texttt{/api/v1/internal/rag/search} and \texttt{/api/v1/internal/rag/generate} using the \texttt{X-Service-Key} internal authentication header.
    \item \textbf{Live screen awareness:} During support calls, the frontend publishes Visual AI screen packets to the backend. Agents can call \texttt{get\_live\_screen\_analysis} or \texttt{generate\_screen\_answer} to answer questions about the user's shared screen.
    \item \textbf{STT/TTS/LLM factory:} In realtime mode, Gemini Realtime can provide an integrated voice model. In pipeline mode, the system combines local \texttt{faster-whisper} STT, Gemini or OpenAI LLMs, and Edge TTS by default, with OpenAI or Google alternatives when configured.
    \item \textbf{Call persistence:} The LiveKit server collects transcripts, records audio tracks into WAV files, translates/finalizes transcripts, and saves a \texttt{voice\_call\_logs} row on shutdown.
    \item \textbf{Escalation:} Agents can call \texttt{escalate\_to\_human\_agent}. The backend creates a high-priority escalated ticket through \texttt{/api/v1/internal/voice/escalations}.
\end{itemize}

\subsection{Voice Call Operations API}

The \texttt{/api/v1/voice-calls} routes expose recorded calls to agents and administrators:

\begin{itemize}
    \item \texttt{GET /voice-calls/}: list call logs with pagination.
    \item \texttt{GET /voice-calls/\{call\_id\}}: inspect a call transcript and metadata.
    \item \texttt{GET /voice-calls/\{call\_id\}/audio}: stream the recorded WAV file.
    \item \texttt{POST /voice-calls/\{call\_id\}/post-call-summary}: generate a structured summary, action items, resolution status, and ticket suggestions.
    \item \texttt{POST /voice-calls/\{call\_id\}/link-ticket}: attach the call to an existing ticket or create a new ticket using the post-call summary.
\end{itemize}

%---------------------------------------------------------------------
\section{Testing and Validation}

Validation covers the Visual AI pipeline, screenshare assistance, LiveKit token generation, voice-agent screen context tools, and post-call workflows.

\begin{table}[H]
\centering
\caption{Sprint 3 Validation Summary}
\label{tab:sprint3-tests}
\renewcommand{\arraystretch}{1.3}
\small
\begin{tabular}{|p{2.4cm}|p{3.0cm}|p{4.1cm}|p{2.5cm}|c|}
\hline
\textbf{Test Case} & \textbf{Preconditions} & \textbf{Expected Result} & \textbf{Postconditions} & \textbf{Status} \\
\hline
Screenshot Analysis & Authenticated user uploads an image & Gemini provider returns OCR, caption, UI labels, and embedding metadata & Analysis can be displayed in the UI & Pass \\
\hline
Gap Detection & Analysis and reference screen exist & Weighted gap score and severity are returned & Guidance hints identify missing or unexpected UI evidence & Pass \\
\hline
Screenshare Frames & User submits ordered frame captures with consent & Low-FPS sampling, transition scoring, final-frame analysis, and assistance hints are returned & Optional support-call context is published & Pass \\
\hline
Realtime Chunk & Browser uploads short WebM/MP4 screen chunks & Frames are extracted and analyzed within the realtime endpoint contract & Voice agents can read latest packet context & Pass \\
\hline
Support Call Token & Authenticated user requests call access & Backend returns a LiveKit room token for \texttt{support-call-\{user\_id\}} & Client can join the support-call room & Pass \\
\hline
Voice-Agent RAG & Caller asks a TunisieSMS knowledge question & Agent calls the internal RAG bridge before answering & Spoken answer is grounded in KB content when available & Pass \\
\hline
Post-Call Workflow & A call log contains transcript or audio metadata & Summary and ticket suggestion are generated & Call can be linked to a ticket & Pass \\
\hline
\end{tabular}
\end{table}

%---------------------------------------------------------------------
\section*{Conclusion}

Sprint 3 delivered a working multimodal support layer. The platform can now analyze screenshots, compare them with reference screens, process live screenshare frames, feed visual context into LiveKit agents, answer voice questions with RAG grounding, persist call transcripts and recordings, and create escalation or follow-up tickets from voice sessions.


%=====================================================================
%=====================================================================

\chapter{Sprint 4 - Frontend, Integration \& Optimization}
\label{chap:sprint4}

\section*{Introduction}

\noindent This chapter describes the final integration sprint. Sprint 4 connects the backend modules to the React frontend, completes the operational views for agents and administrators, containerizes the LiveKit voice-agent runtime, and adds reliability improvements around API keys, runtime settings, notifications, dashboards, and assisted workflows.

\medskip

\textbf{Sprint Goal:} Deliver the integrated user interface, connect Visual AI and LiveKit voice workflows to the frontend, expose operational admin controls, and improve runtime reliability through Docker Compose orchestration, configurable settings, key selection, and end-to-end validation.

%---------------------------------------------------------------------
\section{Sprint Backlog}

\begin{table}[H]
\smaller
\centering
\caption{Sprint 4 Backlog}
\label{tab:sprint4-backlog}
\begin{tabular}{|l|p{4.8cm}|p{4.8cm}|c|}
\hline
\textbf{US \#} & \textbf{Feature} & \textbf{Task} & \textbf{Est.\,(h)} \\
\hline
\multirow{4}{*}{US15} & \multirow{4}{*}{\parbox{4.8cm}{FRONTEND INTERFACES}} & Client ticket, conversation, email, WhatsApp, and notification views & 30 \\
\cline{3-4}
 & & Agent workspace with tickets, conversations, decisions, RAG, and escalations & 30 \\
\cline{3-4}
 & & Admin pages for users, settings, audit logs, chat access, QR bridge, and AI snippets & 30 \\
\cline{3-4}
 & & Visual AI pages: screenshot assistant, screenshare, wizard, and references & 25 \\
\hline
\multirow{3}{*}{US17} & \multirow{3}{*}{\parbox{4.8cm}{VOICE UI}} & Client support-call page with LiveKit connection and screen sharing & 25 \\
\cline{3-4}
 & & Voice call logs, audio playback, post-call summaries, and ticket linking & 20 \\
\cline{3-4}
 & & Voice-agent runtime admin page with start, stop, status, logs, and escalation form & 20 \\
\hline
\multirow{3}{*}{US20} & \multirow{3}{*}{\parbox{4.8cm}{INTEGRATION \& RELIABILITY}} & Docker Compose integration for API, Celery, WhatsApp bridge, LiveKit, and voice agents & 20 \\
\cline{3-4}
 & & Gemini key selection from comma-separated keys for backend and voice agents & 10 \\
\cline{3-4}
 & & Runtime settings, notifications, dashboard caching, and assisted draft workflows & 25 \\
\hline
\multicolumn{3}{|r|}{\textbf{Total}} & \textbf{210} \\
\hline
\end{tabular}
\end{table}

%---------------------------------------------------------------------
\section{Design and Architecture Updates}

Sprint 4 integrates the previously separate modules behind a React single-page application. The frontend uses protected routes and role-based navigation. Clients can open support calls and share screens, agents can handle tickets and recorded calls, and administrators can operate runtime services.

\begin{figure}[H]
\centering
\includegraphics[angle=90,height=0.88\textheight]{class-diagram.png}
\caption{Global Class Diagram - Integrated Platform Architecture}
\label{fig:global-class-diagram}
\end{figure}

%---------------------------------------------------------------------
\section{Implementation}

\subsection{Frontend Delivery}

The frontend is implemented as a Vite React application under \texttt{frontend/}. The application entry point \texttt{frontend/src/App.tsx} wires together \texttt{BrowserRouter}, \texttt{AuthProvider}, \texttt{ProtectedRoute}, React Query, toast notifications, tooltips, lazy-loaded feature pages, and the shared \texttt{AppLayout}. The sidebar in \texttt{frontend/src/shared/components/AppSidebar.tsx} builds role-aware navigation: clients receive dashboard, conversations, tickets, and notifications, while agents and administrators receive the operational support workspace; administrators also receive the administration menu.

\begin{itemize}
    \item \textbf{Authentication and application shell:} \texttt{/landing}, \texttt{/login}, and \texttt{/register} provide public entry points. Authenticated pages are wrapped by \texttt{ProtectedRoute}; role-restricted routes such as \texttt{/voice-agents}, \texttt{/settings}, \texttt{/admins}, \texttt{/agents}, \texttt{/clients}, and \texttt{/logs} are limited to administrators. The shared shell provides the collapsible sidebar, profile shortcut, logout action, language-aware layout direction, loading states, and global API error handling.
    \item \textbf{Dashboard interface:} \texttt{/} renders \texttt{DashboardPage}, a role-aware command dashboard with KPI cards, quick actions, weekly activity, ticket volume and status charts, AI operations metrics, user composition charts, assisted-draft performance, urgent tickets, recent tickets, and client support status. The page consumes dashboard APIs through React Query and presents cached server state with loading and empty states.
    \item \textbf{Ticket interfaces:} \texttt{/tickets} uses separate client and operator views. Clients can search, paginate, create tickets, and open ticket details. Agents and administrators use a ticket command center with filters, date range controls, status and priority badges, assignment actions, and agent selection. \texttt{/tickets/\{id\}} provides operator ticket analysis, status and priority updates, assignment context, and client edit flows. Client ticket details include AI helper actions for rewriting, summarizing, and improving ticket content.
    \item \textbf{Conversation interfaces:} \texttt{/conversations} switches between an operator inbox and a client chat experience. Operators can filter conversations, inspect context panels, manage SLA chips, use AI draft snippets, send replies, and view message attachments. Clients receive a chat-style interface with conversation search, new-chat creation, file attachments, voice-message recording, AI typing feedback, and a direct action to start a LiveKit support call.
    \item \textbf{Email workspace:} \texttt{/email} implements the email support interface with thread grouping, folder and status filters, search, label filtering, compose and reply forms, bulk actions such as read/unread/star updates, body normalization, link rendering, and AI-assisted draft telemetry. The page connects to the email API module and keeps operator actions inside the same support workspace.
    \item \textbf{WhatsApp interfaces:} \texttt{/whatsapp} provides the agent WhatsApp inbox with conversation search, refresh controls, message thread display, reply composition, and AI snippet insertion. \texttt{/supervision} is the administrator supervision view for monitoring WhatsApp agents, searching users and conversations, inspecting selected threads, and copying conversation identifiers. The legacy \texttt{/whatsapp-supervision} route redirects to \texttt{/supervision}.
    \item \textbf{Decision, RAG, and escalation interfaces:} \texttt{/decisions} exposes decision analysis, history, statistics, and playbook tabs. \texttt{/decisions/config} lets administrators tune thresholds, rule toggles, and outcome documentation. \texttt{/rag} provides a knowledge-base workspace with tester, article library, and PDF ingestion tabs. \texttt{/escalations} gives agents and administrators an escalation queue, detailed case overview, timeline, decision evidence, and reply box for escalated tickets.
    \item \textbf{Visual AI interface:} \texttt{/visual-ai} implements the multimodal support workspace. The screenshot tab uploads images and displays OCR, caption, labels, UI elements, confidence, regions, and gap-detection output. The screenshare tab supports captured frame sequences, uploaded videos, realtime chunks, transition scoring, final-frame analysis, assistance hints, and optional reference-key comparison. The troubleshooting wizard tab generates guided troubleshooting steps from the user goal, observed screen state, OCR text, attempted actions, and context hints. Administrators can also manage reference screens with names, keys, descriptions, expected OCR text, expected elements, and deletion actions.
    \item \textbf{Client support-call interface:} \texttt{/support-call} renders the LiveKit client call room. It obtains a room token, joins the support-call room, manages microphone and connection state, shows the voice-agent visual state, starts and stops screen sharing, captures screen frames or realtime recording chunks, sends them to Visual AI, displays live screen-analysis context, and keeps a screen-share preview available during the call.
    \item \textbf{Voice operations interfaces:} \texttt{/voice-calls} lists persisted call logs with room metadata, duration, transcript status, and detail navigation. The detail view plays recorded audio, displays transcripts, generates post-call summaries, suggests action items, links calls to existing tickets, creates new linked tickets, and can escalate unresolved calls. \texttt{/test-agent} provides an agent test room for generating a secure token and connecting to a voice-agent session.
    \item \textbf{Administration interfaces:} \texttt{/admins}, \texttt{/agents}, and \texttt{/clients} reuse the user-management module with role-specific directories, search, status display, create/edit dialogs, and profile metadata. \texttt{/chat-access} manages which users can access chat. \texttt{/qr} provides the WhatsApp bridge setup flow using QR, phone-number, and pairing-code states. \texttt{/voice-agents} exposes the voice-agent runtime dashboard with status, logs, start/stop controls, configuration updates, token generation, and manual escalation tooling. \texttt{/ai-draft-snippets} lets administrators create, search, edit, activate, and document reusable AI response snippets. \texttt{/settings} centralizes application, ticketing, email, SMTP/Gmail, AI, and advanced runtime settings. \texttt{/logs} exposes audit-log search and filters by resource, action, and user.
    \item \textbf{Notification and profile interfaces:} \texttt{/notifications} lists user notifications and supports marking notifications as read. \texttt{/profile} lets users update identity, contact details, timezone, language, avatar URL, and password through guarded profile forms.
    \item \textbf{Shared UI and API layer:} The frontend uses reusable components under \texttt{frontend/src/components/ui} and \texttt{frontend/src/shared/components}, including tables, tabs, dialogs, badges, skeletons, empty states, KPI cards, and status badges. Each feature consumes typed API clients from \texttt{frontend/src/shared/api}, keeping authentication refresh, API error mapping, pagination, and server-state invalidation consistent across all interfaces.
\end{itemize}

\subsection{Containerization \& Runtime Control}

The current Docker Compose file runs PostgreSQL with pgvector, Redis, FastAPI, Celery worker, Celery Beat, Flower, WhatsApp Bridge, LiveKit, and the voice-agent worker.

\begin{itemize}
    \item \textbf{FastAPI service:} Runs Alembic migrations before starting Uvicorn. It supports reload mode for development or multiple workers for normal container execution.
    \item \textbf{Voice-agent service:} Builds from \texttt{voice\_agents/Dockerfile}, shares networking with the LiveKit service, exposes a control API on port \texttt{8001}, and writes recordings to \texttt{/app/recordings}.
    \item \textbf{Control proxy:} Backend routes under \texttt{/voice-agents} call the voice-agent control API to start, stop, check status, and fetch logs.
    \item \textbf{Key selection:} The backend and voice-agent settings support comma-separated \texttt{GEMINI\_API\_KEY} values. The active key is selected randomly through \texttt{current\_gemini\_key}, which distributes requests across available keys.
\end{itemize}

%---------------------------------------------------------------------
\section{Testing and Performance}

Sprint 4 validation focuses on integrated behavior rather than isolated modules only. The repository includes frontend tests for authentication guards, client calls, client chat/tickets, table rendering, and screenshare validation, plus backend tests for assisted drafts, notifications/dashboard caching, voice validation, support-call tokens, screen context, and post-call summaries.

\begin{table}[H]
\centering
\caption{Sprint 4 End-to-End Validation}
\label{tab:sprint4-tests}
\renewcommand{\arraystretch}{1.3}
\small
\begin{tabular}{|p{2.5cm}|p{3.0cm}|p{4.0cm}|p{2.7cm}|c|}
\hline
\textbf{Test Case} & \textbf{Preconditions} & \textbf{Expected Result} & \textbf{Postconditions} & \textbf{Status} \\
\hline
Frontend Routing & User logs in with role-specific account & Protected routes allow only authorized pages & Navigation matches user role & Pass \\
\hline
Visual AI UI & User submits screenshot or screenshare data & Frontend calls Visual AI endpoints and renders analysis & User receives actionable hints & Pass \\
\hline
Support Call & User opens \texttt{/support-call} & LiveKit token is issued and room connection starts & Screen packets can be sent to backend & Pass \\
\hline
Voice Runtime Admin & Admin opens \texttt{/voice-agents} & Status, logs, start, and stop controls call backend proxy & Worker state is visible from UI & Pass \\
\hline
Post-Call Handling & A voice call log exists & Summary and ticket-link actions are available & Ticket references the call through \texttt{source\_voice\_call\_id} & Pass \\
\hline
Docker Composition & All services are started through Compose & API, Redis, PostgreSQL, LiveKit, and voice-agent control API are reachable & Integrated workflow can run locally & Pass \\
\hline
\end{tabular}
\end{table}

\begin{table}[H]
\centering
\caption{Sprint 4 Operational Targets}
\label{tab:sprint4-metrics}
\renewcommand{\arraystretch}{1.2}
\begin{tabular}{|l|c|p{5.2cm}|c|}
\hline
\textbf{Area} & \textbf{Target} & \textbf{Validation Method} & \textbf{Status} \\
\hline
API health & Service responds on \texttt{/health} & FastAPI health endpoint and Docker startup checks & Pass \\
\hline
RAG and AI generation & Provider failures are surfaced as service errors or fallbacks & Backend response provider tests and internal RAG bridge handling & Pass \\
\hline
Screenshare processing & Invalid type, oversized file, and long video are rejected & Visual AI video and realtime chunk endpoint tests & Pass \\
\hline
Voice-agent runtime & Start, stop, status, and logs are available to admins & Voice-agent control API and frontend admin page & Pass \\
\hline
\end{tabular}
\end{table}

%---------------------------------------------------------------------
\section*{Conclusion}

Sprint 4 completes the integration layer around the backend intelligence. The React application now exposes the ticketing, conversation, email, WhatsApp, Visual AI, RAG, voice-call, and administration workflows. Docker Compose coordinates the API, workers, LiveKit, voice agents, and supporting infrastructure. The result is an integrated support platform in which multimodal context, AI assistance, human escalation, and operational monitoring are available from one user-facing system.
